% Options for packages loaded elsewhere
\PassOptionsToPackage{unicode}{hyperref}
\PassOptionsToPackage{hyphens}{url}
\documentclass[
]{article}
\usepackage{xcolor}
\usepackage[margin=1in]{geometry}
\usepackage{amsmath,amssymb}
\setcounter{secnumdepth}{-\maxdimen} % remove section numbering
\usepackage{iftex}
\ifPDFTeX
  \usepackage[T1]{fontenc}
  \usepackage[utf8]{inputenc}
  \usepackage{textcomp} % provide euro and other symbols
\else % if luatex or xetex
  \usepackage{unicode-math} % this also loads fontspec
  \defaultfontfeatures{Scale=MatchLowercase}
  \defaultfontfeatures[\rmfamily]{Ligatures=TeX,Scale=1}
\fi
\usepackage{lmodern}
\ifPDFTeX\else
  % xetex/luatex font selection
\fi
% Use upquote if available, for straight quotes in verbatim environments
\IfFileExists{upquote.sty}{\usepackage{upquote}}{}
\IfFileExists{microtype.sty}{% use microtype if available
  \usepackage[]{microtype}
  \UseMicrotypeSet[protrusion]{basicmath} % disable protrusion for tt fonts
}{}
\makeatletter
\@ifundefined{KOMAClassName}{% if non-KOMA class
  \IfFileExists{parskip.sty}{%
    \usepackage{parskip}
  }{% else
    \setlength{\parindent}{0pt}
    \setlength{\parskip}{6pt plus 2pt minus 1pt}}
}{% if KOMA class
  \KOMAoptions{parskip=half}}
\makeatother
\usepackage{color}
\usepackage{fancyvrb}
\newcommand{\VerbBar}{|}
\newcommand{\VERB}{\Verb[commandchars=\\\{\}]}
\DefineVerbatimEnvironment{Highlighting}{Verbatim}{commandchars=\\\{\}}
% Add ',fontsize=\small' for more characters per line
\usepackage{framed}
\definecolor{shadecolor}{RGB}{248,248,248}
\newenvironment{Shaded}{\begin{snugshade}}{\end{snugshade}}
\newcommand{\AlertTok}[1]{\textcolor[rgb]{0.94,0.16,0.16}{#1}}
\newcommand{\AnnotationTok}[1]{\textcolor[rgb]{0.56,0.35,0.01}{\textbf{\textit{#1}}}}
\newcommand{\AttributeTok}[1]{\textcolor[rgb]{0.13,0.29,0.53}{#1}}
\newcommand{\BaseNTok}[1]{\textcolor[rgb]{0.00,0.00,0.81}{#1}}
\newcommand{\BuiltInTok}[1]{#1}
\newcommand{\CharTok}[1]{\textcolor[rgb]{0.31,0.60,0.02}{#1}}
\newcommand{\CommentTok}[1]{\textcolor[rgb]{0.56,0.35,0.01}{\textit{#1}}}
\newcommand{\CommentVarTok}[1]{\textcolor[rgb]{0.56,0.35,0.01}{\textbf{\textit{#1}}}}
\newcommand{\ConstantTok}[1]{\textcolor[rgb]{0.56,0.35,0.01}{#1}}
\newcommand{\ControlFlowTok}[1]{\textcolor[rgb]{0.13,0.29,0.53}{\textbf{#1}}}
\newcommand{\DataTypeTok}[1]{\textcolor[rgb]{0.13,0.29,0.53}{#1}}
\newcommand{\DecValTok}[1]{\textcolor[rgb]{0.00,0.00,0.81}{#1}}
\newcommand{\DocumentationTok}[1]{\textcolor[rgb]{0.56,0.35,0.01}{\textbf{\textit{#1}}}}
\newcommand{\ErrorTok}[1]{\textcolor[rgb]{0.64,0.00,0.00}{\textbf{#1}}}
\newcommand{\ExtensionTok}[1]{#1}
\newcommand{\FloatTok}[1]{\textcolor[rgb]{0.00,0.00,0.81}{#1}}
\newcommand{\FunctionTok}[1]{\textcolor[rgb]{0.13,0.29,0.53}{\textbf{#1}}}
\newcommand{\ImportTok}[1]{#1}
\newcommand{\InformationTok}[1]{\textcolor[rgb]{0.56,0.35,0.01}{\textbf{\textit{#1}}}}
\newcommand{\KeywordTok}[1]{\textcolor[rgb]{0.13,0.29,0.53}{\textbf{#1}}}
\newcommand{\NormalTok}[1]{#1}
\newcommand{\OperatorTok}[1]{\textcolor[rgb]{0.81,0.36,0.00}{\textbf{#1}}}
\newcommand{\OtherTok}[1]{\textcolor[rgb]{0.56,0.35,0.01}{#1}}
\newcommand{\PreprocessorTok}[1]{\textcolor[rgb]{0.56,0.35,0.01}{\textit{#1}}}
\newcommand{\RegionMarkerTok}[1]{#1}
\newcommand{\SpecialCharTok}[1]{\textcolor[rgb]{0.81,0.36,0.00}{\textbf{#1}}}
\newcommand{\SpecialStringTok}[1]{\textcolor[rgb]{0.31,0.60,0.02}{#1}}
\newcommand{\StringTok}[1]{\textcolor[rgb]{0.31,0.60,0.02}{#1}}
\newcommand{\VariableTok}[1]{\textcolor[rgb]{0.00,0.00,0.00}{#1}}
\newcommand{\VerbatimStringTok}[1]{\textcolor[rgb]{0.31,0.60,0.02}{#1}}
\newcommand{\WarningTok}[1]{\textcolor[rgb]{0.56,0.35,0.01}{\textbf{\textit{#1}}}}
\usepackage{graphicx}
\makeatletter
\newsavebox\pandoc@box
\newcommand*\pandocbounded[1]{% scales image to fit in text height/width
  \sbox\pandoc@box{#1}%
  \Gscale@div\@tempa{\textheight}{\dimexpr\ht\pandoc@box+\dp\pandoc@box\relax}%
  \Gscale@div\@tempb{\linewidth}{\wd\pandoc@box}%
  \ifdim\@tempb\p@<\@tempa\p@\let\@tempa\@tempb\fi% select the smaller of both
  \ifdim\@tempa\p@<\p@\scalebox{\@tempa}{\usebox\pandoc@box}%
  \else\usebox{\pandoc@box}%
  \fi%
}
% Set default figure placement to htbp
\def\fps@figure{htbp}
\makeatother
\setlength{\emergencystretch}{3em} % prevent overfull lines
\providecommand{\tightlist}{%
  \setlength{\itemsep}{0pt}\setlength{\parskip}{0pt}}
\usepackage{bookmark}
\IfFileExists{xurl.sty}{\usepackage{xurl}}{} % add URL line breaks if available
\urlstyle{same}
\hypersetup{
  pdftitle={MANOVA},
  hidelinks,
  pdfcreator={LaTeX via pandoc}}

\title{MANOVA}
\author{}
\date{\vspace{-2.5em}2026-04-17}

\begin{document}
\maketitle

\subsection{Introduction}\label{introduction}

MANOVA --- multivariate analysis of variance --- tests whether group
means differ across several dependent variables simultaneously,
accounting for the correlations among the outcomes. It is more powerful
than running separate one-way ANOVAs when outcomes are correlated,
because it pools information across the response set instead of testing
each in isolation. It also controls Type I error at the family-wise
level for the joint hypothesis, removing the need for
multiple-comparison correction across a battery of univariate ANOVAs.

\subsection{Prerequisites}\label{prerequisites}

A working understanding of one-way ANOVA, multivariate Normal
distributions, and covariance matrices.

\subsection{Theory}\label{theory}

MANOVA decomposes the total cross-product matrix into hypothesis (\(H\))
and error (\(E\)) components and constructs four standard multivariate
test statistics from their joint structure:

\begin{itemize}
\tightlist
\item
  \textbf{Wilks' lambda}: \(\Lambda = |E| / |H + E|\), the classical
  choice; small values reject the null.
\item
  \textbf{Pillai's trace}: \(V = \mathrm{tr}(H (H + E)^{-1})\), robust
  to covariance heterogeneity and the recommended default.
\item
  \textbf{Hotelling--Lawley trace}: \(T = \mathrm{tr}(H E^{-1})\), more
  powerful when one dimension dominates.
\item
  \textbf{Roy's greatest root}: \$\theta = \$ largest eigenvalue of
  \(H E^{-1}\), most powerful when group differences are concentrated in
  one direction but liberal otherwise.
\end{itemize}

All four test the same omnibus null hypothesis but have slightly
different operating characteristics under non-Normality and unequal
covariances.

\subsection{Assumptions}\label{assumptions}

Multivariate Normal within groups, homogeneous covariance matrices
across groups (Box's M test), and independent observations.

\subsection{R Implementation}\label{r-implementation}

\begin{Shaded}
\begin{Highlighting}[]
\NormalTok{fit }\OtherTok{\textless{}{-}} \FunctionTok{manova}\NormalTok{(}\FunctionTok{cbind}\NormalTok{(Sepal.Length, Sepal.Width, Petal.Length, Petal.Width) }\SpecialCharTok{\textasciitilde{}}
\NormalTok{              Species, }\AttributeTok{data =}\NormalTok{ iris)}
\FunctionTok{summary}\NormalTok{(fit)                     }\CommentTok{\# Pillai default}
\FunctionTok{summary}\NormalTok{(fit, }\AttributeTok{test =} \StringTok{"Wilks"}\NormalTok{)}
\FunctionTok{summary}\NormalTok{(fit, }\AttributeTok{test =} \StringTok{"Hotelling{-}Lawley"}\NormalTok{)}

\CommentTok{\# Follow{-}up univariate ANOVAs}
\FunctionTok{summary.aov}\NormalTok{(fit)}

\CommentTok{\# car Anova for Type II/III}
\FunctionTok{library}\NormalTok{(car)}
\FunctionTok{Manova}\NormalTok{(}\FunctionTok{lm}\NormalTok{(}\FunctionTok{cbind}\NormalTok{(Sepal.Length, Sepal.Width) }\SpecialCharTok{\textasciitilde{}}\NormalTok{ Species, }\AttributeTok{data =}\NormalTok{ iris))}
\end{Highlighting}
\end{Shaded}

\subsection{Output \& Results}\label{output-results}

\texttt{summary(fit)} returns the Pillai trace by default, with \(F\)
approximation, degrees of freedom, and \(p\)-value. Specifying other
tests via the \texttt{test} argument produces Wilks, Hotelling-Lawley,
and Roy variants. \texttt{summary.aov(fit)} runs follow-up univariate
ANOVAs on each outcome.

\subsection{Interpretation}\label{interpretation}

A reporting sentence: ``MANOVA showed a significant multivariate species
effect (Wilks' \(\Lambda = 0.02\), \(F(8, 288) = 199\), \(p < 0.001\),
multivariate \(\eta^2 = 0.86\)); univariate follow-up ANOVAs confirmed
significant species differences on each of the four flower dimensions
(\(p < 0.001\) for each), with petal length carrying the strongest
effect (\(\eta^2 = 0.94\)).'' Always pair MANOVA with both univariate
follow-ups and a multivariate effect-size measure.

\subsection{Practical Tips}\label{practical-tips}

\begin{itemize}
\tightlist
\item
  Pillai's trace is the recommended default --- it is the most robust to
  covariance heterogeneity and unequal sample sizes.
\item
  For a significant MANOVA, follow up with univariate ANOVAs on each
  outcome to identify which dimensions drive the effect; corrected
  \(p\)-values (Bonferroni or Holm) are appropriate at this stage.
\item
  Test homogeneity of covariance with Box's M (\texttt{biotools::boxM})
  before relying on Wilks; a significant Box's M motivates Pillai.
\item
  For unequal sample sizes and multifactorial designs, use
  \texttt{car::Manova()} with explicit Type II or Type III sums of
  squares; base R's \texttt{summary.manova} defaults to sequential (Type
  I).
\item
  Discriminant analysis follows naturally as a post-MANOVA exploration;
  the linear discriminants visualise where the multivariate group
  differences live.
\item
  For few groups with many highly-correlated outcomes, MANOVA is more
  powerful than Bonferroni-corrected univariate tests; for many groups
  with few outcomes, separate ANOVAs may be more interpretable.
\end{itemize}

\subsection{R Packages Used}\label{r-packages-used}

Base R \texttt{manova()} and \texttt{summary.manova()};
\texttt{car::Manova()} for Type II / III sums of squares;
\texttt{biotools::boxM()} for the homogeneity-of-covariance test;
\texttt{MASS::lda()} for discriminant-analysis follow-up;
\texttt{heplots} for visualising hypothesis-error decompositions.

\subsection{For Reviewers}\label{for-reviewers}

\textbf{What to look for in a paper using this method.}

\begin{itemize}
\tightlist
\item
  \textbf{Common misapplications.}
\item
  \textbf{Diagnostics that should be reported but often aren't.}
\item
  \textbf{Red flags in tables and figures.}
\item
  \textbf{What to verify.}
\item
  \textbf{What an adequate Methods paragraph must contain.}
\end{itemize}

\subsection{See also --- labs in this
chapter}\label{see-also-labs-in-this-chapter}

\begin{itemize}
\tightlist
\item
  \href{labs/lab-pca-factor-analysis-cca-lda.Rmd}{PCA, factor analysis,
  CCA, LDA}
\item
  \href{labs/lab-clustering-k-means-hierarchical-model-based.Rmd}{Clustering:
  k-means, hierarchical, model-based}
\item
  \href{labs/lab-umap-and-t-sne.Rmd}{UMAP and t-SNE}
\end{itemize}

\end{document}
